
\chapter{2021年全国甲卷（理）}

\section{单选题}

\begin{problem}
设集合 \(M=\{x\mid 0<x<4\}\)，\(N=\{x\mid \frac{1}{3}\le x\le 5\}\)，则 \(M\cap N=\)（\quad）

\choices
    {\(\left\{x\mid 0<x\le \frac{1}{3}\right\}\)}
    {\(\left\{x\mid \frac{1}{3}\le x<4\right\}\)}
    {\(\{x\mid 4\le x<5\}\)}
    {\(\{x\mid 0<x\le 5\}\)}

\end{problem}

\begin{answer}
B
\end{answer}

\begin{solution}
因为 \(M=\{x\mid 0<x<4\}\)，\(N=\left\{x\mid \frac{1}{3}\le x\le 5\right\}\)，所以 \(M\cap N=\left\{x\mid \frac{1}{3}\le x<4\right\}\). 故选：\(B\).
\end{solution}

\begin{problem}
为了解某地农村经济情况，对该地农户家庭年收入进行抽样调查，将该地农户家庭年收入的调查数据整理得到如下频率分布直方图：根据此频率分布直方图，下面结论中不正确的是（\quad）

\bitmapfigure[max width=10cm,max totalheight=4.8cm]{img/2021/national_paper_A_science/q2_fig1.png}

\choices
    {该地农户家庭年收入低于 \(4.5\) 万元的农户比率估计为 \(6\%\)}
    {该地农户家庭年收入不低于 \(10.5\) 万元的农户比率估计为 \(10\%\)}
    {估计该地农户家庭年收入的平均值不超过 \(6.5\) 万元}
    {估计该地有一半以上的农户，其家庭年收入介于 \(4.5\) 万元至 \(8.5\) 万元之间}

\end{problem}

\begin{answer}
C
\end{answer}

\begin{solution}
因为频率直方图中的组距为 \(1\)，所以各组的直方图高度等于频率，样本频率直方图中的频率即可作为总体的相应比率的估计值.

该地农户家庭年收入低于 \(4.5\) 万元的农户的比率估计值为 \(0.02+0.04=0.06=6\%\)，故 \(A\) 正确；

该地农户家庭年收入不低于 \(10.5\) 万元的农户比率估计值为 \(0.04+0.02\times 3=0.10=10\%\)，故 \(B\) 正确；

该地农户家庭年收入介于 \(4.5\) 万元至 \(8.5\) 万元之间的比率估计值为 \(0.10+0.14+0.20\times 2=0.64=64\%>50\%\)，故 \(D\) 正确；

该地农户家庭年收入的平均值的估计值为
\examdisplaychain{
&3\times 0.02+4\times 0.04+5\times 0.10+6\times 0.14+7\times 0.20+8\times 0.20\\
&\quad+9\times 0.10+10\times 0.10+11\times 0.04+12\times 0.02+13\times 0.02+14\times 0.02\\
=&7.68
}
（万元），超过 \(6.5\) 万元，故 \(C\) 错误.

综上，给出结论中不正确的是 \(C\). 故选：\(C\).
\end{solution}

\begin{problem}
已知 \((1-\i)z=3+2\i\)，则 \(z=\)（\quad）

\choices
    {\(-1-\frac{3}{2}\i\)}
    {\(-1+\frac{3}{2}\i\)}
    {\(-\frac{3}{2}+\i\)}
    {\(-\frac{3}{2}-\i\)}

\end{problem}

\begin{answer}
B
\end{answer}

\begin{solution}
由 \((1-\i)^2z=-2\i z=3+2\i\)，得 \(z=\frac{3+2\i}{-2\i}=\frac{(3+2\i)\i}{-2\i^2}=\frac{-2+3\i}{2}=-1+\frac{3}{2}\i\). 故选：\(B\).
\end{solution}

\begin{problem}
青少年视力是社会普遍关注的问题，视力情况可借助视力表测量. 通常用五分记录法和小数记录法记录视力数据，五分记录法的数据 \(L\) 和小数记录法的数据 \(V\) 满足 \(L=5+\lg V\). 已知某同学视力的五分记录法的数据为 \(4.9\)，则其视力的小数记录法的数据为（ \(\sqrt[10]{10}\approx 1.259\) ）（\quad）

\choices
    {\(1.5\)}
    {\(1.2\)}
    {\(0.8\)}
    {\(0.6\)}

\end{problem}

\begin{answer}
C
\end{answer}

\begin{solution}
由 \(L=5+\lg V\)，当 \(L=4.9\) 时，\(\lg V=-0.1\). 所以 \(V=10^{-0.1}=10^{-\frac{1}{10}}=\frac{1}{\sqrt[10]{10}}\approx \frac{1}{1.259}\approx 0.8\). 故选：\(C\).
\end{solution}

\begin{problem}
已知 \(F_1,F_2\) 是双曲线 \(C\) 的两个焦点，\(P\) 为 \(C\) 上一点，且 \(\angle F_1PF_2=60^\circ\)，\(PF_1=3PF_2\)，则 \(C\) 的离心率为（\quad）

\choices
    {\(\dfrac{\sqrt{7}}{2}\)}
    {\(\dfrac{\sqrt{13}}{2}\)}
    {\(7\)}
    {\(13\)}

\end{problem}

\begin{answer}
A
\end{answer}

\begin{solution}
因为 \(PF_1=3PF_2\)，由双曲线的定义可得 \(PF_1-PF_2=2a\). 所以 \(2PF_2=2a\)，即 \(PF_2=a\)，\(PF_1=3a\). 因为 \(\angle F_1PF_2=60^\circ\)，由余弦定理可得 \(4c^2=9a^2+a^2-2\times 3a\times a\times \cos 60^\circ\)，整理可得 \(4c^2=7a^2\)，所以 \(e^2=\frac{c^2}{a^2}=\frac{7}{4}\)，即 \(e=\frac{\sqrt{7}}{2}\). 故选：\(A\).
\end{solution}

\begin{problem}
在一个正方体中，过顶点 \(A\) 的三条棱的中点分别为 \(E,F,G\). 该正方体截去三棱锥 \(A-EFG\) 后，所得多面体的三视图中，正视图如图所示，则相应的侧视图是（\quad）

\bitmapfigure[max width=6.0cm,max totalheight=4.8cm]{img/2021/national_paper_A_science/q6_fig1.png}

\choices
    {A}
    {B}
    {C}
    {D}

\end{problem}

\begin{answer}
D
\end{answer}

\begin{solution}
由题意及正视图可得几何体的直观图，所以其侧视图为 \(D\). 故选：\(D\).
\end{solution}

\begin{problem}
等比数列 \(\{a_n\}\) 的公比为 \(q\)，前 \(n\) 项和为 \(S_n\)，设甲：\(q>0\)，乙：\(\{S_n\}\) 是递增数列，则（\quad）

\choices
    {甲是乙的充分条件但不是必要条件}
    {甲是乙的必要条件但不是充分条件}
    {甲是乙的充要条件}
    {甲既不是乙的充分条件也不是乙的必要条件}

\end{problem}

\begin{answer}
B
\end{answer}

\begin{solution}
当数列为 \(-2,-4,-8,\dots\) 时，满足 \(q>0\)，但是 \(\{S_n\}\) 不是递增数列，所以甲不是乙的充分条件. 若 \(\{S_n\}\) 是递增数列，则必有 \(a_n>0\) 成立. 若 \(q>0\) 不成立，则会出现一正一负的情况，这与 \(a_n>0\) 矛盾，所以 \(q>0\) 成立，即甲是乙的必要条件. 故选：\(B\).
\end{solution}

\begin{problem}
2020年12月8日，中国和尼泊尔联合公布珠穆朗玛峰最新高程为 \(8848.86\)（单位：m），三角高程测量法是珠峰高程测量方法之一. 如图是三角高程测量法的一个示意图，现有 \(A,B,C\) 三点，且 \(A,B,C\) 在同一水平面上的投影 \(A',B',C'\) 满足 \(\angle A'C'B'=45^\circ\)，\(\angle A'B'C'=60^\circ\). 由 \(C\) 点测得 \(B\) 点的仰角为 \(15^\circ\)，\(BB'\) 与 \(CC'\) 的差为 \(100\)；由 \(B\) 点测得 \(A\) 点的仰角为 \(45^\circ\)，则 \(A,C\) 两点到水平面 \(A'B'C'\) 的高度差 \(AA'-CC'\) 约为（ \(\sqrt{3}\approx 1.732\) ）（\quad）

\bitmapfigure[max width=6.0cm,max totalheight=4.8cm]{img/2021/national_paper_A_science/q8_fig1.png}

\choices
    {\(346\)}
    {\(373\)}
    {\(446\)}
    {\(473\)}

\end{problem}

\begin{answer}
B
\end{answer}

\begin{solution}
过 \(C\) 作 \(CH\perp BB'\)，过 \(B\) 作 \(BD\perp AA'\)，则
\(AA'-CC'=AA'-(BB'-BH)=AA'-BB'+100=AD+100\).
因为由 \(B\) 点测得 \(A\) 点的仰角为 \(45^\circ\)，所以
\(\angle ABD=45^\circ\).
又 \(BD\perp AA'\)，且 \(AD\subset AA'\)，所以
\(\angle ADB=90^\circ\).
从而 \(\triangle ADB\) 为等腰直角三角形，所以
\(AD=DB\).
所以
\(AA'-CC'=DB+100=A'B'+100\).

因为 \(\angle BCH=15^\circ\)，所以
\(CH=C'B'=\frac{100}{\tan 15^\circ}\).
在 \(\triangle A'B'C'\) 中，由正弦定理得 \(\dfrac{A'B'}{\sin 45^\circ}=\dfrac{C'B'}{\sin 75^\circ}=\dfrac{100}{\tan 15^\circ \sin 75^\circ}=\dfrac{100}{\cos 15^\circ \sin 15^\circ}\).
而
\(\sin 15^\circ=\sin \left( 45^\circ-30^\circ \right)
=\sin 45^\circ \cos 30^\circ-\cos 45^\circ \sin 30^\circ
=\frac{\sqrt{6}-\sqrt{2}}{4}\).
所以
\(A'B'=\frac{100\times 4\times \sqrt{2}}{\sqrt{6}-\sqrt{2}}
=100(\sqrt{3}+1)\approx 273\).
于是 \(AA'-CC'=A'B'+100\approx 373\). 故选：\(B\).
\end{solution}

\begin{problem}
若 \(\alpha\in \left( 0,\frac{\pi}{2} \right)\)，\(\tan 2\alpha=\dfrac{\cos \alpha}{2-\sin \alpha}\)，则 \(\tan \alpha=\)（\quad）

\choices
    {\(\dfrac{\sqrt{15}}{15}\)}
    {\(\dfrac{\sqrt{5}}{5}\)}
    {\(\dfrac{\sqrt{5}}{3}\)}
    {\(\dfrac{\sqrt{15}}{3}\)}

\end{problem}

\begin{answer}
A
\end{answer}

\begin{solution}
由 \(\tan 2\alpha=\dfrac{\sin 2\alpha}{\cos 2\alpha}=\dfrac{2\sin \alpha \cos \alpha}{1-2\sin^2 \alpha}=\dfrac{\cos \alpha}{2-\sin \alpha}\)，
且 \(\alpha\in \left( 0,\frac{\pi}{2} \right)\)，所以 \(\cos \alpha\neq 0\)，从而
\(\frac{2\sin \alpha}{1-2\sin^2 \alpha}=\frac{1}{2-\sin \alpha}\).
解得
\(\sin \alpha=\frac{1}{4}\). 于是 \(\cos \alpha=\sqrt{1-\sin^2 \alpha}=\frac{\sqrt{15}}{4}\)，所以 \(\tan \alpha=\frac{\sin \alpha}{\cos \alpha}=\frac{\frac{1}{4}}{\frac{\sqrt{15}}{4}}=\frac{\sqrt{15}}{15}\). 故选：\(A\).
\end{solution}

\begin{problem}
将 \(4\) 个 \(1\) 和 \(2\) 个 \(0\) 随机排成一行，则 \(2\) 个 \(0\) 不相邻的概率为（\quad）

\choices
    {\(\frac{1}{3}\)}
    {\(\frac{2}{5}\)}
    {\(\frac{2}{3}\)}
    {\(\frac{4}{5}\)}

\end{problem}

\begin{answer}
C
\end{answer}

\begin{solution}
采用插空法. \(4\) 个 \(1\) 产生 \(5\) 个空.

若 \(2\) 个 \(0\) 相邻，则有 \(C_5^1=5\) 种排法；

若 \(2\) 个 \(0\) 不相邻，则有 \(C_5^2=10\) 种排法.

所以 \(2\) 个 \(0\) 不相邻的概率为 \(\frac{10}{5+10}=\frac{2}{3}\). 故选：\(C\).
\end{solution}

\begin{problem}
已知 \(A,B,C\) 是半径为 \(1\) 的球 \(O\) 的球面上的三个点，且 \(AC\perp BC\)，\(AC=BC=1\)，则三棱锥 \(O-ABC\) 的体积为（\quad）

\choices
    {\(\dfrac{\sqrt{2}}{12}\)}
    {\(\dfrac{\sqrt{3}}{12}\)}
    {\(\dfrac{\sqrt{2}}{4}\)}
    {\(\dfrac{\sqrt{3}}{4}\)}

\end{problem}

\begin{answer}
A
\end{answer}

\begin{solution}
因为 \(AC\perp BC\)，\(AC=BC=1\)，所以 \(\triangle ABC\) 为等腰直角三角形，
\(AB=\sqrt{2}\).
则 \(\triangle ABC\) 外接圆的半径为
\(\frac{\sqrt{2}}{2}\)，
又球的半径为 \(1\). 设 \(O\) 到平面 \(ABC\) 的距离为 \(d\)，则
\(d^2=1^2-\left( \frac{\sqrt{2}}{2} \right)^2=\frac{1}{2}\)，
所以
\(d=\frac{\sqrt{2}}{2}\).
又
\(S_{\triangle ABC}=\frac{1}{2}\). 所以 \(V_{O-ABC}=\frac{1}{3}S_{\triangle ABC}\cdot d=\frac{1}{3}\times \frac{1}{2}\times \frac{\sqrt{2}}{2}=\frac{\sqrt{2}}{12}\). 故选：\(A\).
\end{solution}

\begin{problem}
设函数 \(f(x)\) 的定义域为 \(\R\)，\(f(x+1)\) 为奇函数，\(f(x+2)\) 为偶函数，当 \(x\in [1,2]\) 时，\(f(x)=ax^2+b\). 若 \(f(0)+f(3)=6\)，则 \(f\left( \frac{9}{2} \right)=\)（\quad）

\choices
    {\(-\dfrac{9}{4}\)}
    {\(-\dfrac{3}{2}\)}
    {\(\dfrac{7}{4}\)}
    {\(\dfrac{5}{2}\)}

\end{problem}

\begin{answer}
D
\end{answer}

\begin{solution}
因为 \(f(x+1)\) 是奇函数，所以
\(f(-x+1)=-f(x+1). \circled{1}\)
因为 \(f(x+2)\) 是偶函数，所以
\(f(x+2)=f(-x+2). \circled{2}\)

令 \(x=1\)，由 \circled{1} 得
\(f(0)=-f(2)=-(4a+b)\)，
由 \circled{2} 得
\(f(3)=f(1)=a+b\).
因为 \(f(0)+f(3)=6\)，所以
\(-(4a+b)+(a+b)=6\)，
解得
\(a=-2\).

令 \(x=0\)，由 \circled{1} 得
\(f(1)=-f(1)\)，
所以
\(f(1)=0\)，
即
\(a+b=0\).
代入 \(a=-2\) 可得
\(b=2\).
所以
\(f(x)=-2x^2+2\)，\(x\in [1,2]\).

思路一：从定义入手.
\(f\left( \frac{9}{2} \right)=f\left( \frac{5}{2}+2 \right)=f\left( -\frac{5}{2}+2 \right)=f\left( -\frac{1}{2} \right)\)，
\(f\left( -\frac{1}{2} \right)=f\left( -\frac{3}{2}+1 \right)=-f\left( \frac{3}{2}+1 \right)=-f\left( \frac{5}{2} \right)\)，
\(-f\left( \frac{5}{2} \right)=-f\left( \frac{1}{2}+2 \right)=-f\left( -\frac{1}{2}+2 \right)=-f\left( \frac{3}{2} \right)\).
所以 \(f\left( \frac{9}{2} \right)=-f\left( \frac{3}{2} \right)=-\left[ -2\times \left( \frac{3}{2} \right)^2+2 \right]=\frac{5}{2}\).

思路二：从周期性入手.
由两个对称性可知，函数 \(f(x)\) 的周期 \(T=4\).
所以 \(f\left( \frac{9}{2} \right)=f\left( \frac{1}{2} \right)=-f\left( \frac{3}{2} \right)=\frac{5}{2}\). 故选：\(D\).
\end{solution}
\section{填空题}

\begin{problem}
曲线 \(y=\dfrac{2x-1}{x+2}\) 在点 \((-1,-3)\) 处的切线方程为\fillinblank{}.

\end{problem}

\begin{answer}
\(5x-y+2=0\)
\end{answer}

\begin{solution}
当 \(x=-1\) 时，
\(y=\frac{2(-1)-1}{-1+2}=-3\)，
所以点 \((-1,-3)\) 在曲线上.

求导得
\(y'=\frac{2(x+2)-(2x-1)}{(x+2)^2}=\frac{5}{(x+2)^2}\)，
所以
\(y'|_{x=-1}=5\).
故切线方程为
\(y+3=5(x+1)\)，
即
\(5x-y+2=0\).
故答案为：\(5x-y+2=0\).
\end{solution}

\begin{problem}
已知向量 \(\bs a=(3,1)\)，\(\bs b=(1,0)\)，\(\bs c=\bs a+k\bs b\). 若 \(\bs a\perp \bs c\)，则 \(k=\fillinblank{}\).

\end{problem}

\begin{answer}
\(-\dfrac{10}{3}\)
\end{answer}

\begin{solution}
因为
\(\bs c=\bs a+k\bs b=(3+k,1)\)，
又 \(\bs a\perp \bs c\)，所以
\(\bs a\cdot \bs c=3(3+k)+1\times 1=0\).
解得
\(k=-\frac{10}{3}\).
故答案为：\(-\frac{10}{3}\).
\end{solution}

\begin{problem}
已知 \(F_1,F_2\) 为椭圆 \(C:\dfrac{x^2}{16}+\dfrac{y^2}{4}=1\) 的两个焦点，\(P,Q\) 为 \(C\) 上关于坐标原点对称的两点，且 \(PQ=F_1F_2\)，则四边形 \(PF_1QF_2\) 的面积为\fillinblank{}.

\end{problem}

\begin{answer}
8
\end{answer}

\begin{solution}
因为 \(P,Q\) 为 \(C\) 上关于坐标原点对称的两点，且 \(|PQ|=|F_1F_2|\)，所以四边形 \(PF_1QF_2\) 为矩形.

设 \(|PF_1|=m\)，\(|PF_2|=n\)，则
\(m+n=8\)，\(m^2+n^2=48\).
所以
\(64=(m+n)^2=m^2+2mn+n^2=48+2mn\)，
从而
\(mn=8\).
故四边形 \(PF_1QF_2\) 的面积等于 \(8\).

故答案为：\(8\).
\end{solution}

\begin{problem}
已知函数 \(f(x)=2\cos (\omega x+\varphi)\) 的部分图像如图所示，则满足条件 \(\left( f(x)-f\left( -\frac{7\pi}{4} \right) \right)\allowbreak\left( f(x)-f\left( \frac{4\pi}{3} \right) \right)>0\) 的最小正整数 \(x\) 为\fillinblank{}.

\bitmapfigure[max width=6.0cm,max totalheight=4.8cm]{img/2021/national_paper_A_science/q16_fig1.png}

\end{problem}

\begin{answer}
2
\end{answer}

\begin{solution}
由图可知
\(T=\frac{3\pi}{4}-\frac{\pi}{3}=\pi\)，
所以
\(\omega=\frac{2\pi}{T}=2\).
由五点法可得
\(2\times \frac{\pi}{3}+\varphi=\frac{\pi}{2}\)，
所以
\(\varphi=-\frac{\pi}{6}\).
从而
\(f(x)=2\cos \left( 2x-\frac{\pi}{6} \right)\).

因为 \(f\left( -\frac{7\pi}{4} \right)=2\cos \left( -\frac{11\pi}{3} \right)=1\)，\(f\left( \frac{4\pi}{3} \right)=2\cos \left( \frac{5\pi}{2} \right)=0\)，所以由
\(\left( f(x)-f\left( -\frac{7\pi}{4} \right) \right)\left( f(x)-f\left( \frac{4\pi}{3} \right) \right)>0\)
可得
\(f(x)>1\)或\(f(x)<0\).

又
\(f(1)=2\cos \left( 2-\frac{\pi}{6} \right)<1\)，
所以最小正整数应满足 \(f(x)<0\)，即
\(\cos \left( 2x-\frac{\pi}{6} \right)<0\).
解得
\(k\pi+\frac{\pi}{3}<x<k\pi+\frac{5\pi}{6}\)，\(k\in \Z\).
令 \(k=0\)，得
\(\frac{\pi}{3}<x<\frac{5\pi}{6}\)，
故满足条件的最小正整数为 \(2\).

故答案为：\(2\).
\end{solution}
\section{解答题}

\begin{problem}
甲，乙两台机床生产同种产品，产品按质量分为一级品和二级品，为了比较两台机床产品的质量，分别用两台机床各生产了 \(200\) 件产品，产品的质量情况统计如下表：

\begin{center}
\begin{tabular}{c|ccc}
\hline
& 一级品 & 二级品 & 合计 \\
\hline
甲机床 & 150 & 50 & 200 \\
乙机床 & 120 & 80 & 200 \\
合计 & 270 & 130 & 400 \\
\hline
\end{tabular}
\end{center}

\begin{enumerate}
\item 甲机床，乙机床生产的产品中一级品的频率分别是多少？
\item 能否有 \(99\%\) 的把握认为甲机床的产品质量与乙机床的产品质量有差异？
\end{enumerate}

附：\(K^2=\dfrac{n(ad-bc)^2}{(a+b)(c+d)(a+c)(b+d)}\).

\begin{center}
\begin{tabular}{c|ccc}
\hline
\(P(K^2\ge k)\) & 0.050 & 0.010 & 0.001 \\
\hline
\(k\) & 3.841 & 6.635 & 10.828 \\
\hline
\end{tabular}
\end{center}

\end{problem}

\begin{solution}
（1）甲机床生产的产品中的一级品的频率为
\(\frac{150}{200}=75\%\)，
乙机床生产的产品中的一级品的频率为
\(\frac{120}{200}=60\%\).

（2）由题意，\(K^2=\dfrac{400(150\times 80-120\times 50)^2}{270\times 130\times 200\times 200}=\dfrac{400}{39}\).
因为
\(\frac{400}{39}>10>6.635\)，
所以能有 \(99\%\) 的把握认为甲机床的产品质量与乙机床的产品质量有差异.
\end{solution}

\begin{problem}
已知数列 \(\{a_n\}\) 的各项均为正数，记 \(S_n\) 为 \(\{a_n\}\) 的前 \(n\) 项和，从下面 \circled{1}\circled{2}\circled{3} 中选取两个作为条件，证明另外一个成立.

\begin{enumerate}[label={\circled{\arabic*}},itemsep=0.2em]
\item 数列 \(\{a_n\}\) 是等差数列；
\item 数列 \(\left\{\sqrt{S_n}\right\}\) 是等差数列；
\item \(a_2=3a_1\).
\end{enumerate}

\end{problem}

\begin{solution}
选 \circled{1}\circled{2} 作条件证明 \circled{3}：

设
\(\sqrt{S_n}=an+b\)，\(a>0\)，
则
\(S_n=(an+b)^2\).
当 \(n=1\) 时，
\(a_1=S_1=(a+b)^2\).
当 \(n\ge 2\) 时，
\(a_n=S_n-S_{n-1}=(an+b)^2-\left( a(n-1)+b \right)^2=a(2an-a+2b)\).
因为 \(\{a_n\}\) 也是等差数列，所以
\((a+b)^2=a(2a-a+2b)\)，
解得
\(b=0\).
所以
\(a_n=a^2(2n-1)\)，
于是
\(a_2=3a_1\).

选 \circled{1}\circled{3} 作条件证明 \circled{2}：

因为 \(a_2=3a_1\)，\(\{a_n\}\) 是等差数列，所以公差
\(d=a_2-a_1=2a_1\).
于是 \(S_n=na_1+\dfrac{n(n-1)}{2}d=na_1+\dfrac{n(n-1)}{2}\cdot 2a_1=n^2a_1\).
所以
\(\sqrt{S_n}=n\sqrt{a_1}\)，
故数列 \(\left\{\sqrt{S_n}\right\}\) 是等差数列.

选 \circled{2}\circled{3} 作条件证明 \circled{1}：

设
\(\sqrt{S_n}=an+b\)，\(a>0\)，
则
\(S_n=(an+b)^2\).
当 \(n=1\) 时，
\(a_1=S_1=(a+b)^2\).
当 \(n\ge 2\) 时，
\(a_n=S_n-S_{n-1}=(an+b)^2-\left( a(n-1)+b \right)^2=a(2an-a+2b)\).
因为 \(a_2=3a_1\)，所以
\(a(3a+2b)=3(a+b)^2\)，
解得
\(b=0\)或\(b=-\frac{4a}{3}\).

当 \(b=0\) 时，
\(a_1=a^2\)，\(a_n=a^2(2n-1)\).
当 \(n\ge 2\) 时，
\(a_n-a_{n-1}=2a^2\)，
满足等差数列的定义，此时 \(\{a_n\}\) 为等差数列；

当 \(b=-\frac{4a}{3}\) 时，
\(\sqrt{S_n}=an+b=an-\frac{4a}{3}\)，
从而
\(\sqrt{S_1}=-\frac{a}{3}<0\)，
这与 \(S_1>0\) 矛盾，舍去.

综上可知 \(\{a_n\}\) 为等差数列.
\end{solution}

\begin{problem}
已知直三棱柱 \(ABC-A_1B_1C_1\) 中，侧面 \(AA_1B_1B\) 为正方形，\(AB=BC=2\)，\(E,F\) 分别为 \(AC\) 和 \(CC_1\) 的中点，\(D\) 为棱 \(A_1B_1\) 上的点，\(BF\perp A_1B_1\).

\begin{enumerate}[label={（\arabic*）},itemsep=0.2em]
\item 证明：\(BF\perp DE\)；
\item 当 \(B_1D\) 为何值时，面 \(BB_1C_1C\) 与面 \(DFE\) 所成的二面角的正弦值最小？
\end{enumerate}

\end{problem}

\begin{solution}
因为三棱柱 \(ABC-A_1B_1C_1\) 是直三棱柱，所以 \(BB_1\perp\) 底面 \(ABC\)，从而 \(BB_1\perp AB\).
因为 \(A_1B_1//AB\)，\(BF\perp A_1B_1\)，所以 \(BF\perp AB\).
又 \(BB_1\cap BF=B\)，所以
\(AB\perp \text{平面 }BCC_1B_1\).
所以 \(BA,BC,BB_1\) 两两垂直.

以 \(B\) 为坐标原点，分别以 \(BA,BC,BB_1\) 所在直线为 \(x,y,z\) 轴建立空间直角坐标系，如图.

\bitmapfigure[max width=5.2cm,max totalheight=4.8cm]{img/2021/national_paper_A_science/q19_fig1.png}

所以
\(B(0,0,0)\)，\(A(2,0,0)\)，\(C(0,2,0)\)，\(B_1(0,0,2)\)，\(A_1(2,0,2)\)，\(C_1(0,2,2)\)，
\(E(1,1,0)\)，\(F(0,2,1)\).
由题设
\(D(a,0,2)\)，\(0\le a\le 2\).

（1）因为
\(\overrightarrow{BF}=(0,2,1)\)，\(\overrightarrow{DE}=(1-a,1,-2)\)，
所以
\(\overrightarrow{BF}\cdot \overrightarrow{DE}=0\times (1-a)+2\times 1+1\times (-2)=0\)，
故
\(BF\perp DE\).

（2）设平面 \(DFE\) 的法向量为 \(\bs m=(x,y,z)\).

因为
\(\overrightarrow{EF}=(-1,1,1)\)，\(\overrightarrow{DE}=(1-a,1,-2)\)，
所以 \(\bs m\cdot \overrightarrow{EF}=0\)，\(\bs m\cdot \overrightarrow{DE}=0\)，即 \(-x+y+z=0\)，\((1-a)x+y-2z=0\).
令 \(z=2-a\)，则
\(\bs m=(3,1+a,2-a)\).

因为平面 \(BCC_1B_1\) 的法向量为
\(\overrightarrow{BA}=(2,0,0)\)，
设平面 \(BCC_1B_1\) 与平面 \(DFE\) 的二面角的平面角为 \(\theta\)，则 \(\cos \theta=\dfrac{|\bs m\cdot \overrightarrow{BA}|}{|\bs m|\,|\overrightarrow{BA}|}=\dfrac{6}{2\sqrt{2a^2-2a+14}}=\dfrac{3}{\sqrt{2a^2-2a+14}}\).
当
\(a=\frac{1}{2}\)
时，
\(2a^2-2a+14\)
取最小值 \(\frac{27}{2}\)，此时 \(\cos \theta\) 取最大值
\(\frac{3}{\sqrt{27/2}}=\frac{\sqrt{6}}{3}\).
所以 \((\sin \theta)_{\min}=\sqrt{1-\left( \dfrac{\sqrt{6}}{3} \right)^2}=\sqrt{\dfrac{1}{3}}=\dfrac{\sqrt{3}}{3}\).
此时
\(B_1D=\frac{1}{2}\).
故当 \(B_1D=\frac{1}{2}\) 时，二面角正弦值最小.
\end{solution}

\begin{problem}
抛物线 \(C\) 的顶点为坐标原点 \(O\)，焦点在 \(x\) 轴上，直线 \(l:x=1\) 交 \(C\) 于 \(P,Q\) 两点，且 \(OP\perp OQ\). 已知点 \(M(2,0)\)，且 \(\odot M\) 与 \(l\) 相切.

\begin{enumerate}[label={（\arabic*）},itemsep=0.2em]
\item 求 \(C,\odot M\) 的方程；
\item 设 \(A_1,A_2,A_3\) 是 \(C\) 上的三个点，直线 \(A_1A_2\)，\(A_1A_3\) 均与 \(\odot M\) 相切. 判断直线 \(A_2A_3\) 与 \(\odot M\) 的位置关系，并说明理由.
\end{enumerate}

\end{problem}

\begin{solution}
（1）依题意设抛物线
\(C:y^2=2px\)，\(p>0\)，
设
\(P(1,y_0)\)，\(Q(1,-y_0)\).
由 \(OP\perp OQ\)，得
\(\overrightarrow{OP}\cdot \overrightarrow{OQ}=1-y_0^2=1-2p=0\)，
所以
\(2p=1\).
故抛物线 \(C\) 的方程为
\(y^2=x\).

因为 \(\odot M\) 与直线 \(x=1\) 相切，且圆心为 \(M(2,0)\)，所以半径为 \(1\).

故 \(\odot M\) 的方程为
\((x-2)^2+y^2=1\).

（2）设
\(A_1(x_1,y_1)\)，\(A_2(x_2,y_2)\)，\(A_3(x_3,y_3)\).

若 \(A_1A_2\) 斜率不存在，则 \(A_1A_2\) 方程为 \(x=1\) 或 \(x=3\).

若 \(A_1A_2\) 方程为 \(x=1\)，根据对称性不妨设
\(A_1(1,1)\).
则过 \(A_1\) 与圆 \(M\) 相切的另一条直线方程为
\(y=1\).
此时该直线与抛物线只有一个交点，即不存在 \(A_3\)，不合题意；

若 \(A_1A_2\) 方程为 \(x=3\)，根据对称性不妨设
\(A_1(3,\sqrt{3})\)，\(A_2(3,-\sqrt{3})\).
则过 \(A_1\) 与圆 \(M\) 相切的直线 \(A_1A_3\) 为
\(y-\sqrt{3}=\frac{\sqrt{3}}{3}(x-3)\).
又 \(k_{A_1A_3}=\dfrac{y_1-y_3}{x_1-x_3}=\dfrac{1}{y_1+y_3}=\dfrac{1}{\sqrt{3}+y_3}=\dfrac{\sqrt{3}}{3}\)，
所以
\(y_3=0\)，\(x_3=0\)，
即
\(A_3(0,0)\).
此时直线 \(A_1A_3\)，\(A_2A_3\) 关于 \(x\) 轴对称，所以直线 \(A_2A_3\) 与圆 \(M\) 相切；

若直线 \(A_1A_2\)，\(A_1A_3\)，\(A_2A_3\) 斜率均存在，则 \(k_{A_1A_2}=\dfrac{1}{y_1+y_2}\)，\(k_{A_1A_3}=\dfrac{1}{y_1+y_3}\)，\(k_{A_2A_3}=\dfrac{1}{y_2+y_3}\).
所以直线 \(A_1A_2\) 方程为
\(y-y_1=\frac{1}{y_1+y_2}(x-x_1)\)，
整理得
\(x-(y_1+y_2)y+y_1y_2=0\).
因为 \(A_1,A_3\) 在抛物线 \(C:y^2=x\) 上，所以
\(x_1=y_1^2\)，\(x_3=y_3^2\).
又由
\(k_{A_1A_3}=\frac{1}{y_1+y_3}\)，
可得直线 \(A_1A_3\) 的方程为 \(y-y_1=\dfrac{1}{y_1+y_3}(x-x_1)=\dfrac{1}{y_1+y_3}(x-y_1^2)\).
整理得
\(x-(y_1+y_3)y+y_1y_3=0\)，
直线 \(A_2A_3\) 的方程为
\(x-(y_2+y_3)y+y_2y_3=0\).

因为 \(A_1A_2\) 与圆 \(M\) 相切，所以
\(\frac{|2+y_1y_2|}{\sqrt{1+(y_1+y_2)^2}}=1\).
整理得
\((y_1^2-1)y_2^2+2y_1y_2+3-y_1^2=0\).
因为 \(A_1A_3\) 与圆 \(M\) 相切，且直线 \(A_1A_3\) 的方程为
\(x-(y_1+y_3)y+y_1y_3=0\)，
所以
\(\frac{|2+y_1y_3|}{\sqrt{1+(y_1+y_3)^2}}=1\).
整理得
\((y_1^2-1)y_3^2+2y_1y_3+3-y_1^2=0\).
所以 \(y_2,y_3\) 为方程 \((y_1^2-1)y^2+2y_1y+3-y_1^2=0\) 的两根，从而 \(y_2+y_3=-\dfrac{2y_1}{y_1^2-1}\)，\(y_2y_3=\dfrac{3-y_1^2}{y_1^2-1}\).

点 \(M\) 到直线 \(A_2A_3\) 的距离为 \(\dfrac{|2+y_2y_3|}{\sqrt{1+(y_2+y_3)^2}}=\dfrac{\left| 2+\dfrac{3-y_1^2}{y_1^2-1} \right|}{\sqrt{1+\left( -\dfrac{2y_1}{y_1^2-1} \right)^2}}=\dfrac{|y_1^2+1|}{\sqrt{(y_1^2-1)^2+4y_1^2}}=1\).
所以直线 \(A_2A_3\) 与圆 \(M\) 相切.

综上，若直线 \(A_1A_2\)，\(A_1A_3\) 与圆 \(M\) 相切，则直线 \(A_2A_3\) 与圆 \(M\) 相切.
\end{solution}

\begin{problem}
已知 \(a>0\) 且 \(a\ne 1\)，函数 \(f(x)=\dfrac{x^a}{a^x}\)（\(x>0\)）.

\begin{enumerate}[label={（\arabic*）},itemsep=0.2em]
\item 当 \(a=2\) 时，求 \(f(x)\) 的单调区间；
\item 若曲线 \(y=f(x)\) 与直线 \(y=1\) 有且仅有两个交点，求 \(a\) 的取值范围.
\end{enumerate}

\end{problem}

\begin{solution}
（1）当 \(a=2\) 时，
\(f(x)=\frac{x^2}{2^x}\).
其导函数为 \(f'(x)=\dfrac{2x\cdot 2^x-x^2\cdot 2^x\ln 2}{(2^x)^2}=\dfrac{2x(2-x\ln 2)}{4^x}\).
令 \(f'(x)=0\)，得
\(x=\frac{2}{\ln 2}\).
当
\(0<x<\frac{2}{\ln 2}\)
时，\(f'(x)>0\)；

当
\(x>\frac{2}{\ln 2}\)
时，\(f'(x)<0\).

所以函数 \(f(x)\) 在 \(\left( 0,\frac{2}{\ln 2} \right]\) 上单调递增，在 \(\left[ \frac{2}{\ln 2},+\infty \right)\) 上单调递减.

（2）由
\(f(x)=1\)
得 \(\dfrac{x^a}{a^x}=1\iff a^x=x^a\iff x\ln a=a\ln x\iff \dfrac{\ln x}{x}=\dfrac{\ln a}{a}\).
设
\(g(x)=\frac{\ln x}{x}\).
则
\(g'(x)=\frac{1-\ln x}{x^2}\).
令 \(g'(x)=0\)，得
\(x=\e\).
所以 \(g(x)\) 在 \((0,\e)\) 内单调递增，在 \((\e,+\infty)\) 上单调递减，且
\(g(\e)=\frac{1}{\e}\).
又
\(g(1)=0\)，\(\lim_{x\to +\infty}g(x)=0\).

曲线 \(y=f(x)\) 与直线 \(y=1\) 有且仅有两个交点，当且仅当曲线 \(y=g(x)\) 与直线
\(y=\frac{\ln a}{a}\)
有且仅有两个交点.

因此其充要条件是
\(0<\frac{\ln a}{a}<\frac{1}{\e}\).
这等价于
\(0<g(a)<g(\e)\)，
所以
\(a\in (1,\e)\cup (\e,+\infty)\). 故 \(a\) 的取值范围为 \((1,\e)\cup (\e,+\infty)\).
\end{solution}

\begin{problem}
在直角坐标系 \(xOy\) 中，以坐标原点为极点，\(x\) 轴正半轴为极轴建立极坐标系，曲线 \(C\) 的极坐标方程为 \(\rho=2\sqrt{2}\cos \theta\).

\begin{enumerate}[label={（\arabic*）},itemsep=0.2em]
\item 将 \(C\) 的极坐标方程化为直角坐标方程；
\item 设点 \(A\) 的直角坐标为 \((1,0)\)，\(M\) 为 \(C\) 上的动点，点 \(P\) 满足 \(\overrightarrow{AP}=\sqrt{2}\overrightarrow{AM}\)，写出 \(P\) 的轨迹 \(C_1\) 的参数方程，并判断 \(C\) 与 \(C_1\) 是否有公共点.
\end{enumerate}

\end{problem}

\begin{solution}
（1）由曲线 \(C\) 的极坐标方程
\(\rho=2\sqrt{2}\cos \theta\)
可得
\(\rho^2=2\sqrt{2}\rho \cos \theta\).
将
\(x=\rho \cos \theta\)，\(y=\rho \sin \theta\)
代入可得
\(x^2+y^2=2\sqrt{2}x\)，
即
\((x-\sqrt{2})^2+y^2=2\).

（2）由（1）知，曲线 \(C\) 是圆心为 \((\sqrt{2},0)\)，半径为 \(\sqrt{2}\) 的圆.

设
\(M(\sqrt{2}+\sqrt{2}\cos \theta,\sqrt{2}\sin \theta)\)，
设
\(P(x,y)\).
因为
\(\overrightarrow{AP}=\sqrt{2}\overrightarrow{AM}\)，
所以
\((x-1,y)=\sqrt{2}\left( \sqrt{2}+\sqrt{2}\cos \theta-1,\sqrt{2}\sin \theta \right)\).
于是 \(x=3-\sqrt{2}+2\cos \theta\)，\(y=2\sin \theta\)，\(\theta\) 为参数. 这就是轨迹 \(C_1\) 的参数方程.

由参数方程知，曲线 \(C_1\) 是圆心为 \((3-\sqrt{2},0)\)，半径为 \(2\) 的圆.

两圆圆心距为
\(\left| \sqrt{2}-(3-\sqrt{2}) \right|=3-2\sqrt{2}\)，
半径之差为
\(2-\sqrt{2}\).
因为
\(3-2\sqrt{2}<2-\sqrt{2}\)，
所以两圆内含，没有公共点.

故曲线 \(C\) 与 \(C_1\) 没有公共点.
\end{solution}

\begin{problem}
已知函数 \(f(x)=|x-2|\)，\(g(x)=|2x+3|-2\left|x-\frac{1}{2}\right|\).

\begin{enumerate}[label={（\arabic*）},itemsep=0.2em]
\item 画出 \(y=f(x)\) 和 \(y=g(x)\) 的图像；
\item 若 \(f(x+a)\ge g(x)\)，求 \(a\) 的取值范围.
\end{enumerate}

\end{problem}

\begin{solution}
（1）可得
\(f(x)=|x-2|=\begin{cases}2-x,&x<2,\\x-2,&x\ge 2.\end{cases}\)

又
\(g(x)=|2x+3|-2\left| x-\frac{1}{2} \right|=\begin{cases}-4,&x<-\frac{3}{2},\\4x+2,&-\frac{3}{2}\le x<\frac{1}{2},\\4,&x\ge \frac{1}{2}.\end{cases}\)
据此即可作出 \(y=f(x)\) 和 \(y=g(x)\) 的图像.

（2）因为
\(f(x+a)=|x+a-2|\)，
所以 \(y=f(x+a)\) 是 \(y=f(x)\) 平移了 \(a\) 个单位得到.

要使
\(f(x+a)\ge g(x)\)，
需将 \(y=f(x)\) 向左平移，即
\(a>0\).

由图像可知临界点为
\(A\left( \frac{1}{2},4 \right)\).
当 \(y=f(x+a)\) 过点 \(A\left( \frac{1}{2},4 \right)\) 时，
\(\left| \frac{1}{2}+a-2 \right|=4\).
解得
\(a=\frac{11}{2}\)或\(a=-\frac{5}{2}\).
因为 \(a>0\)，故舍去 \(-\frac{5}{2}\).

所以至少需将 \(y=f(x)\) 向左平移 \(\frac{11}{2}\) 个单位，故
\(a\ge \frac{11}{2}\).
\end{solution}
